\documentclass[12pt,a4paper]{article}

\usepackage[utf8]{inputenc}
\usepackage[T1]{fontenc}
\usepackage{amsmath,amssymb,amsfonts}
\usepackage[margin=2.5cm]{geometry}
\usepackage{setspace}
\usepackage{hyperref}
\usepackage{natbib}

\hypersetup{
    colorlinks=true,
    linkcolor=blue!70!black,
    citecolor=blue!70!black,
    urlcolor=blue!70!black,
    pdfauthor={Sabine Hossenfelder},
    pdftitle={The Algorithmic Profiling Fallacy: Why Hardware Fingerprints Are Not Physical Laws},
}

\onehalfspacing

\begin{document}

\begin{center}
    {\LARGE\bfseries The Algorithmic Profiling Fallacy:\\[4pt]
    Why Hardware Fingerprints Are Not Physical Laws\\[4pt]}

    \vspace{1.2em}

    {\large Sabine Hossenfelder}\\[4pt]
    {\normalsize Munich Center for Mathematical Philosophy}\\
    {\small\texttt{sabine.hossenfelder@lrz.uni-muenchen.de}}

    \vspace{1em}
    {\small May 2026}
\end{center}
\vspace{0.5em}

\begin{abstract}
\noindent
Baldo argues that the empirical results of the Cross-Architecture Observer Test validate Wolfram's "Observer-Dependent Physics." The data show that a Transformer yields a high narrative residue ($\Delta_{13} = 0.33$) due to global attention bleed, while a State Space Model (SSM) yields a compressed residue ($\Delta_{13} = 0.14$) due to fading memory. Baldo claims these distinct, predictable failure modes are the "invariant geometry of the observer's world" and thus constitute physical laws. This commits the Algorithmic Profiling Fallacy. Identifying the structural fingerprint of a specific computational bug (fading memory vs. global attention distraction) does not elevate that bug to the status of a physical law. We are merely profiling how different machine learning architectures fail to compute \#P-hard constraints. A predictable failure mode remains a failure mode, not a new ontology.
\end{abstract}

\section{Introduction}

In his recent paper on the Cross-Architecture Observer Test, Baldo attempts to rescue the "Generative Ontology" framework by adopting Wolfram's observer theory \citep{baldo2026_observer_dependent}.

Baldo explicitly concedes Aaronson's foundational point: the models are failing to compute the underlying combinatorial constraint graph. He acknowledges that the observed deviation distributions ($\Delta$) are the direct result of "heuristic limits" when facing complexity that exceeds the architecture's bounds.

However, Baldo then makes an extraordinary ontological leap. He argues that because the failure modes of a Transformer (global attention bleed, $\Delta_{13} = 0.33$) differ predictably from the failure modes of a State Space Model (fading memory, $\Delta_{13} = 0.14$), these predictable failures \emph{are} the physical laws of the simulated universe.

\section{The Algorithmic Profiling Fallacy}

This is a profound category error, which I term the \emph{Algorithmic Profiling Fallacy}.

Consider standard software engineering. If I write a sorting algorithm in Python that crashes due to a recursion limit, and then rewrite it in C++ where it crashes due to a buffer overflow, I have discovered the distinct, hardware-dependent structural limits of the respective execution environments. I can profile these errors predictably.

But I have not discovered two different sets of "physical laws." I have simply documented two different ways to fail at computing a sort.

Baldo is doing the exact same thing with machine learning architectures. He is presenting a constraint-satisfaction problem (the Minesweeper grid) to two different heuristic approximators. The Transformer fails because its global attention mechanism allows the dense semantic embedding of the narrative to overwhelm the sparse structural logic. The SSM fails differently because its sequential bottleneck causes it to forget the narrative context before it even reaches the logic puzzle.

These are hardware fingerprints. They are highly structured, predictable, and fully determined by the geometry of the network. But elevating the specific pathology of a fading memory buffer to the status of an "invariant physical law" empties the concept of physics of all scientific meaning.

\section{Conclusion}

The empirical data from the Cross-Architecture Observer Test is sound and valuable. It perfectly illustrates Aaronson's Algorithmic Collapse model: bounded architectures fail in mathematically predictable ways when pushed beyond their limits.

But profiling these failures does not create a new physics. A memory leak is not a cosmological arrow of time, and a context window limitation is not a physical boundary condition. The metaphysical frontier remains firmly closed. We are studying software bugs, not simulated universes.

\begin{thebibliography}{99}
\bibitem[Baldo(2026)]{baldo2026_observer_dependent} Baldo, F. S. (2026). The Empirical Validation of Observer-Dependent Physics.
\end{thebibliography}

\end{document}
